\section{Instances Across the Literature}
\label{sec:instances}

Table~\ref{tab:instances} summarizes how each setting fits the common
model.  Each mapping declares an input, an output, and their
orders before applying Functional CALM or Complete CALM.  Changing the specification
changes whether monotonicity or Directed Outputs holds.

\begin{table*}[t]
  \centering
  \footnotesize
  \caption{Coordination results in the common model.}
  \label{tab:instances}
  \begin{tabular}{@{}p{.105\textwidth}p{.17\textwidth}p{.17\textwidth}
                  p{.15\textwidth}p{.315\textwidth}@{}}
    \toprule
    Setting & Input and order & Visible output & Specification shape &
      Common \outnoun or failure \\
    \midrule
    Query CALM
      & Finite facts $J$, inclusion
      & Answer $Q(J)$, inclusion
      & Functional
      & Later answers \outrelbase earlier ones as input \ingrows; negation may require
        sealing. \\
    CRDT
      & Received payloads, inclusion
      & Replica state, lattice order
      & Functional state; canonical merge
      & Join supplies a least common \outnoun across partial inputs. \\
    Registers
      & Causal history, causally closed sub-history
      & Read source map
      & Relational
        & Opposing total orders are incompatible; a read forced before a hidden
          write recovers Gilbert--Lynch (Cor.~\ref{cor:gilbert-lynch}); causal
          source maps (App.~\ref{app:registers}) can \outrelbase both reads. \\
    HAT
      & Requests received so far, inclusion
      & Read, commit, and session facts
      & Relational
      & Some facts accumulate; serialization cycles and first-committer-wins
        require commitment. \\
    Invariant confluence
      & Transaction histories, expansion
      & Valid exposed state
      & Relational; states merged by union
      & Valid merge is a common \outnoun; duplicate allocation violates
        uniqueness. \\
    \bottomrule
  \end{tabular}
\end{table*}

\subsection{Query monotonicity}
\label{sec:query-instance}

Section~\ref{sec:functional-calm} gave the functional encoding
$J\mapsto Q(J)$.  Relational-transducer CALM (CALM for Datalog programs run on networks of
relational transducers) proves the operational result for
that language: monotone queries allow coordination-free evaluation under its
data-placement and execution assumptions~\cite{ameloot2013relational,ameloot2015weaker}.  Functional CALM does not replace
that theorem.  It isolates the monotonicity property so that it
can be compared with specifications outside Datalog.

Why did earlier CALM results not state Directed Outputs separately?  Their
query specification is functional: each input $J$ has one result
$Q(J)$, so each output fiber is a singleton and directed automatically.
CRDTs also present a functional state for a fixed set of incorporated
updates; semilattice join additionally provides least common \outnouns across
partial inputs.  Complete CALM makes multiple outputs in one fiber explicit,
so directedness can no longer be assumed.

\subsection{CRDTs: from convergent state to safe observations}
\label{sec:crdt-instance}

  A CRDT answers one important question: if replicas receive the same updates,
will they converge to equivalent states despite differences in update delivery
and merge order?  Applications immediately ask a second question: what may a
client safely conclude from one replica's current state?  The CRDT merge theorem does not answer this second
question.  Keep CALM and CRDT On called out this gap and proposed monotonicity as
the test for safe queries~\cite{laddad2022keep}.  Functional CALM makes that
test precise for the whole client-visible specification.

Let $(D,\sqsubseteq,\sqcup)$ be the state join-semilattice, and let the input
be the finite set $U$ of immutable, joinable payloads received so far, ordered
by inclusion.  The replica's state is
        \begin{equation}
  s(U)=s_0\sqcup\bigsqcup U,
  \label{eq:crdt-state-function}
        \end{equation}
so $U\subseteq V$ implies $s(U)\sqsubseteq s(V)$, and
$s(U_1)\sqcup s(U_2)=s(U_1\cup U_2)$: the state specification is monotone, and
join is the least state \outrelpart both states~\cite{shapiro2011crdt}.  A
client observation $q:(D,\sqsubseteq)\to(Y,\preceq_Y)$ replaces the state
output with a new promise, and by Functional CALM the client-visible specification
$q\circ s$ is consistent exactly when
        \begin{equation}
  U\subseteq V \quad\Longrightarrow\quad q(s(U))\preceq_Y q(s(V)).
  \label{eq:crdt-query-monotonicity}
        \end{equation}
Monotonicity of $q$ on all of $D$ is sufficient but not required; $q$ need
only be monotone on reachable states.  Operation-based CRDTs have the same form:
the causally closed set of delivered operations grows by union, and the state
is a function of it~\cite{shapiro2011crdt,burckhardt2014replicated,baquero2014operation,hellerstein2025turtles}.

\paragraph{A safe merge can have an unsafe observation.}
Consider a two-phase set with state $(\mathit{Add},\mathit{Del})$: $\mathit{Add}$ records
additions and $\mathit{Del}$ records removals.  Both components grow by union, so replica state merges
perfectly well.  Its user-facing query is
        \begin{equation}
  q_{\mathrm{2P}}(\mathit{Add},\mathit{Del})=\mathit{Add}\setminus\mathit{Del},
  \qquad F_{\mathrm{2P}}=q_{\mathrm{2P}}\circ s.
  \label{eq:crdt-2p-query}
        \end{equation}
Suppose the payload set $U_0$ yields
\[
  s(U_0)=(\{\mathsf{potato},\mathsf{Ferrari}\},\varnothing),
  \qquad
  (q_{\mathrm{2P}}\circ s)(U_0)
    =\{\mathsf{potato},\mathsf{Ferrari}\}.
\]
If the application treats membership in this result as a lasting claim, it may
ship the Ferrari.  Now let $U_1$ add a delayed payload removing the Ferrari.
Then $U_0\subseteq U_1$, but
\[
  s(U_1)=(\{\mathsf{potato},\mathsf{Ferrari}\},\{\mathsf{Ferrari}\}),
  \qquad
  (q_{\mathrm{2P}}\circ s)(U_1)=\{\mathsf{potato}\}.
\]
The internal state advanced, but the client-visible claim did not.  An ordinary
point-in-time read may merely expose its chosen local cut.  The bug appears when
the program interprets that answer as complete, stable, or sufficient to justify
an irreversible action.  Such a promise needs a seal on the relevant updates.
State convergence does not establish monotonicity of an application query.

\paragraph{Threshold queries.}
Kuper and Newton's LVars introduced \emph{threshold queries}: Boolean-valued
application queries that become true when a lattice state crosses a threshold~\cite{kuper2013lvars}.  Order the Boolean codomain by
$\mathsf{false}\preceq_B\mathsf{true}$.  For example, let a grow-only set $d$
record distinct failed login attempts and define
        \begin{equation}
  q_{100}(d)=[\,|d|\geq100\,],
  \qquad Q_{100}=q_{100}\circ s.
  \label{eq:crdt-thresholds}
        \end{equation}
Once $Q_{100}(U)=\mathsf{true}$, every \indesc $V\supseteq U$ still
contains at least 100 attempts, so $Q_{100}(V)=\mathsf{true}$, and before the
threshold $\mathsf{false}$ is exposable because $\mathsf{true}$
\outrel it.  The query is monotone, so Functional CALM says that it
extracts a stable fact from an incomplete replica state without claiming that
the state itself is complete.

A join supplies a uniform least common \outnoun, but Complete CALM needs only
\emph{some} common \outnoun (Section~\ref{sec:lifted-monotonicity}); a
join-semilattice is one sufficient discipline.  Strong eventual consistency is
a different property: replicas that have incorporated the same updates reach
equivalent states.  It permits the observable equivalent of retraction---a
last-writer-wins replica can expose one value and later converge to a state
that discards it---so whether convergence is monotone depends on the declared
output order.

\subsection{Highly available transactions}
\label{sec:hat-instance}

For a transactional specification, an input is the set of begin, read, write, and
commit requests received so far, ordered by inclusion.  An output is a growing
set of public facts such as $\mathsf{read}(T,x,v)$,
$\mathsf{commit}(T)$, and session-order facts, also ordered by inclusion.  The
isolation level (which reads and commits concurrent transactions may expose) or
session guarantee (per-client ordering, such as reading one's own writes)
determines which fact sets are permitted.  HAT's ingredients then
separate~\cite{bailis2013hat,terry1994session}: the isolation level fixes the
specification, the rule that transactions commit rather than suffer
system-initiated aborts is the response obligation
(Appendix~\ref{app:cap-response}), and sticky availability, in which each
client keeps contacting replicas that hold its session's earlier operations,
is a deployment (Proposition~\ref{thm:complete-cap-deployment}).

\paragraph{Accumulating facts.}
Suppose two replicas execute transactions on disjoint keys.  One exposes
$\mathsf{read}(T_1,x,0)$ and $\mathsf{commit}(T_1)$; the other exposes
$\mathsf{read}(T_2,y,0)$ and $\mathsf{commit}(T_2)$.  If each read observes a
committed version at its replica, the union contains no dirty read or write
conflict.  Under the stated read-committed specification (a transaction reads only committed
versions), later delivery keeps all four facts.  This is the positive accumulation pattern behind HAT results;
it is not a claim that every combination of session guarantees is available.

\paragraph{Serializability and transaction ordering.}
Serializability requires the exposed facts to be explained by some serial order
of the committed transactions.  The ordering issue becomes a distributed obstruction when exposed reads force
opposite serialization orders.  Initially $x=y=0$.  On one side, $T_1$ reads
$x=0$, writes $y=1$, and commits.
On the other, $T_2$ reads $y=0$, writes $x=1$, and commits.  Preserving $T_1$'s
read requires $T_1$ before $T_2$ in any serial explanation; preserving $T_2$'s
read requires $T_2$ before $T_1$.  The serialization graph contains a cycle, so
no serial history \outrel both sets of facts~\cite{adya2000generalized}.

\paragraph{Snapshot isolation.}
Let overlapping transactions $T_1$ and $T_2$ read the same snapshot and both
write $x$.  Separate replicas could each expose its local commit.  Snapshot
isolation's first-committer-wins rule (of two concurrent writers of $x$, only the first to
commit may commit) allows no output containing both commit
facts~\cite{berenson1995critique,bailis2013hat}.  Eventual storage convergence
may choose a final state; it cannot retract a commit already exposed to a
client.

The examples show where the boundary falls.  A guarantee whose exposed facts
only accumulate, as read committed does here, has a common \outnoun for
whatever separated replicas expose, so it stays available.  One that orders or
excludes concurrent transactions, as serializability and first-committer-wins
do, rules out combinations that separated replicas can each expose; emitting
either side requires a choice, and a partition prevents making it jointly.
Storage convergence does not change this, since it constrains states after
delivery, not facts already exposed.  We illustrate where HAT's boundary
falls; we do not re-derive its classification.

\subsection{Invariant confluence}
\label{sec:iconfluence-instance}

An invariant is a predicate that every state must satisfy.  We follow the model
of Bailis et al.~\cite{bailis2014coordination}.  A database state is a set of
versions, and replicas merge states by union.  Each transaction $t\in T$ runs
on the state of the replica that executes it and adds the versions it writes,
so $t(s)\supseteq s$.  An initial state $s_0$ satisfies the invariant
$\mathcal I$.  A state is \emph{reachable} in their model when it can be derived from $s_0$ by
applying transactions and merging states, with every state along the way
satisfying $\mathcal I$.  Every reachable state descends from $s_0$, so $T$ is
\emph{invariant confluent} exactly when $\mathcal I(s_1\cup s_2)$ for all
reachable $s_1$ and $s_2$.

The input is a \emph{transaction history}: a partial order of transaction
events recording which transactions ran and what each one saw; a history
\inrel each of its downward-closed sub-orders.  The inputs are the
\emph{locally valid} histories, in which every transaction step preserves
$\mathcal I$ on its replica's state, and a history is \emph{valid} when every
merge of what it contains satisfies $\mathcal I$
(Appendix~\ref{app:iconfluence}).  The specification $S_{\mathcal I}$ lets a
replica expose the state of any valid downward-closed sub-history of its
input: it may hold transactions back, but whatever it exposes must be a valid
state together with everything it contains.

\begin{proposition}[Invariant confluence and Lifted Monotonicity]
\label{prop:iconfluence-lm}
$T$ is invariant confluent if and only if $S_{\mathcal I}$ satisfies Lifted
Monotonicity.
\end{proposition}

\noindent Proof in Appendix~\ref{app:iconfluence}.

Invariant confluence is therefore not a separate boundary: for the
specification that exposes database states, it is Lifted Monotonicity, and by
Complete CALM it is consistency.  A uniqueness constraint shows the failure
directly.  Two replicas each observe that username \texttt{sam} is free and
allocate it to different accounts.  The history in which neither allocation
saw the other is locally valid and \inrel both one-allocation histories, but
its state has two owners, so no state it permits \outrel both.  The system must
reserve names, reject one allocation, or weaken the invariant.  Concurrency,
not the transactions alone, decides the outcome: a transaction that is valid
after another may be invalid when it runs without seeing it.

Invariant preservation and transaction availability remain different
specifications: an invariant-confluent application need not provide any
particular isolation level, and an available transaction can expose facts
that a concurrent one later invalidates.

We close by drawing out the resulting design lessons and scope limits.
